\section{Policy experiment and welfare measures}
\label{ch:question}

Any welfare comparison is meaningful only relative to a stated policy and
baseline. The analysis therefore begins by defining the marginal
infrastructure change, then distinguishes the traffic-based and
spatial-equilibrium measures and states which benefits they omit.

\subsection{The policy experiment}

Suppose an analyst can lower the underlying transport cost on one network
link. Which link produces the largest economic benefit? A traditional answer
multiplies the cost reduction by the traffic using that link. This is the
social-savings calculation associated with \citet{fogel1962quantitative} and
\citet{Fogel1964Railroads}. It is simple, data-light, and often informative.
It does not require the analyst to model changes in routes, modes, prices, or
locations. Such adjustment does not by itself invalidate social savings. The
relevant question is whether those responses change exposure to congestion or
spatial externalities that give reallocation a first-order welfare value.

The package evaluates a local version of that question. A policy is a small
change in the primitive cost of a declared edge--mode. For a directed edge
\(e=(k,l)\) and mode \(m\), define the primitive log cost as
\[
  \theta_{klm}=\log\bar\kappa_{kl,m}.
\]
Here \(\bar\kappa_{kl,m}\) is the multiplicative transport or generalized
commuting cost supplied by infrastructure before endogenous congestion. The
policy derivative is
\[
  -\frac{d\log W}{d\theta_{klm}}.
\]
Here \(W\) is aggregate welfare in the spatial model defined in Section
\ref{ch:equilibrium}. The reported quantity is the elasticity of that welfare
measure with respect to a reduction in primitive edge--mode cost. Multiplying
it by 0.01 gives the first-order welfare change from a one-percent improvement.
The software reports the elasticity and records the shock size used for this
conversion in the run manifest.

The distinction between a directed edge and a physical link matters. A
directed policy changes one direction. A physical-link policy changes the two
reciprocal directions simultaneously and sums their first-order effects. The
package does not treat an unpaired directed edge as bidirectional.

The package evaluates this experiment in either the economic-geography or
urban-commuting model. Section \ref{ch:equilibrium} defines their distinct
economic adjustments and the transport equations they share. The welfare
measures below apply to both models, with traffic expressed in the units
appropriate to each.

\subsection{Traditional and extended measures}

Let \(\Xi_{kl,m}\) be traffic on the edge--mode in the normalization of the
selected spatial model. It is a world-income share in the
economic-geography model and a share of the commuter flow in the urban model.
The Traditional approach is
\[
  H_{klm}=\Xi_{kl,m}.
\]
It is a welfare elasticity, not a raw vehicle count. In the
economic-geography model, multiplying by world income converts the traffic
share to currency units before applying a small proportional cost change. In
the urban model, the conversion depends on the generalized commuting-cost
units used to construct the baseline.

Two full-model effects add the responses omitted from this traffic measure.
The full realized-cost effect
changes the cost users face after congestion. The Extended approach reports
the full primitive-cost effect of the infrastructure-side improvement.
Pass-through separates these experiments because induced traffic can attenuate
the user-facing cost reduction. Both effects value the interaction of route,
modal, congestion, and spatial responses with pre-existing wedges. In the
urban model, the spatial responses are residence and workplace choices
rather than trade, wages, and mobile labor.

The economic-geography formula developed in Section \ref{ch:ift} can be
written as four factors:
\[
\boxed{
\begin{aligned}
\text{Economic-geography Extended benefit}
={}&
\text{traffic}
\times
\text{primitive-to-user pass-through}\\
&{}\times
\text{common spatial scale}
\times
\text{endpoint access value}.
\end{aligned}}
\]
Traffic gives the direct scale of the affected activity. Pass-through measures
the share of the primitive improvement that reaches users after transport
feedback.
The spatial scale depends on the common externality parameters. The endpoint
multipliers measure the marginal welfare value of improving access at the
link's origin and destination after the rest of the spatial economy adjusts.

\begin{figure}[H]
\centering
\begin{tikzpicture}[
  node distance=0.36cm and 0.36cm,
  every node/.style={font=\small,align=center},
  stage/.style={draw=guideblue!45,fill=guideblue!4,rounded corners=2pt,
    minimum height=0.72cm,text width=1.92cm},
  flow/.style={-{Latex[length=2mm]},line width=0.7pt,draw=black!62},
  feedback/.style={-{Latex[length=2mm]},line width=0.7pt,draw=guideorange,
    bend left=34}
]
\node[stage] (primitive) {primitive edge--mode cost};
\node[stage,right=of primitive] (realized) {realized user-facing cost};
\node[stage,right=of realized] (route) {modes and recursive routes};
\node[stage,right=of route] (spatial) {trade or commuting equilibrium};
\node[stage,right=of spatial] (welfare) {aggregate welfare};
\draw[flow] (primitive) -- node[above,font=\scriptsize]{pass-through} (realized);
\draw[flow] (realized) -- (route);
\draw[flow] (route) -- node[above,font=\scriptsize]{traffic} (spatial);
\draw[flow] (spatial) -- (welfare);
\draw[feedback] (spatial.south) to node[below,font=\scriptsize,text=guideorange]
  {traffic and congestion feedback} (realized.south);
\end{tikzpicture}
\caption{Policy transmission and welfare. The policy changes a primitive infrastructure
cost. Congestion and modal adjustment determine the realized cost users face;
routes transmit that change to trade or commuting flows; the selected spatial
model maps the resulting allocation into welfare.}
\label{fig:guide-mechanism}
\end{figure}

Local analysis separates the model's economic response from assumptions about
project scale and permits a common ranking of many candidate links. The
adjoint method also computes all link derivatives from one equilibrium
linearization. These advantages come with a clear limit: a large project can
leave the neighborhood in which the derivative is accurate. Section
\ref{ch:diagnostics} explains how to compare the first-order prediction with
finite changes.

\subsection{Benefits included in the welfare measure}

The welfare measure covers only effects represented in the selected model:
resource or generalized costs of movement, route and mode substitution,
congestion, spatial externalities, and changes in economic location.
Construction, maintenance, financing, safety, environmental, neighborhood,
displacement, segregation, and distributional effects remain separate entries
in a complete appraisal.

\warningbox{Do not label the output ``net social benefit'' unless all
material project costs and externalities have been brought into the welfare
criterion. In the paper application, ``model-implied economic benefit'' is the
accurate description.}

The package therefore complements conventional transport appraisal rather
than replacing it. Conventional analysis supplies engineering costs, demand
evidence, safety effects, and local externalities. The spatial statistic shows
the contribution of observed network use and equilibrium adjustment to the
benefit side of that ledger. For conventional appraisal, see
\citet{SmallVerhoefLindsey2024} and \citet{dePalmaEtAl2011Handbook}; for
spatial effects, see \citet{ReddingTurner2015} and
\citet{donaldson2025transport}.

\subsection{Using the elasticity in project appraisal}

Let \(E_p=-d\log W/d\theta_p\) be the reported benefit elasticity and let
\(s_p>0\) be a small proportional primitive-cost reduction. The local welfare
change is
\[
 \Delta\log W \simeq E_p s_p.
\]
This statement is invariant to the currency used in the input files. A
money-metric conversion requires an additional normalization. If \(W\) is
measured in units for which the chosen income aggregate \(Y_0\) is the relevant
baseline expenditure function, a first-order equivalent-variation
approximation is
\[
 B_0 \simeq E_p s_p Y_0.
\]
For economic geography, the paper normalizes value traffic by world income,
so the corresponding baseline aggregate is \(Y_0=Y^W\), measured over the
same period as the income and traffic data. For an urban application, the
commuter-share normalization alone does not supply a currency conversion; the
analyst must specify the generalized-cost units and the money-metric welfare
normalization.

The package reports \(E_p\), not an annualized project value. Treat \(B_0\) as
an annual flow only when the baseline income, traffic, and welfare
normalization all refer to that annual period. A conventional appraisal then
adds omitted benefits and subtracts construction, maintenance, financing, and
other costs at their own dates:
\[
 \operatorname{NPV}
 =
 \sum_t
 \frac{B_t+B_t^{\mathrm{other}}-C_t}{(1+r)^t}.
\]
Transfers such as toll revenue require separate treatment from resource costs.
Large projects also require a nonlinear counterfactual or other evidence that
the local approximation remains accurate at the proposed scale.

The quick start in Section \ref{ch:quickstart} implements these requirements as a
data-to-results workflow. The detailed argument now turns to the two spatial
models and the transport block they share. Section \ref{ch:welfare} then
establishes when traffic alone is exact and derives the correction when
wedges matter, before Section \ref{ch:implementation} translates the formulas
and required inputs into data and code.
